\scanpage{4}

\uline{Modules des courbes elliptiques}

Soit $C$ une ``courbe elliptique'' homogène sur $S$ loc. noethérien. La section nulle $s$ définit un bon diviseur $D$ sur $C$ \inserted{transversal aux fibres}, d'où un \deleted{fibré}, \inserted{faisceau inv. $\underline L$} dont les \inserted{germes de} sections sont des \deleted{fonctions} germes de fonctions sur $C$ (précisément celles qui sont des ``multiples'' de $D$). Noter que \deleted{$s^*(\underline L)\simeq$}, $C$ étant un schéma en groupes
\[
\unsure{\underline T_{C/S}|\underline{\mathcal T}_{C/S}}=f^*(s^*(\underline T_{C/S}))=f^*(\mathcal L_{C/S})
\]
\margin{[$\mathcal L$ : faisceau de Lie]}

\deleted{\uline{Proposition 1} les $L^n$ $(n\in\mathbb Z)$ sont \unsure{cohomolgt|cohomologt} plats en \unsure{ttes|toute} dimension.}
\[
s^*(\underline L)\simeq s^*(\underline L\underset C\otimes\underline O_{S'})\simeq s^*(\underline N_{C/S'})\simeq s^*(\underline T_{C/S})\simeq\unsure{\mathcal L}
\]
Ce faisceau sera noté $\mathcal L_{C/S}$ [``Lie'']. Donc
\[
s^*(\underline L^{\otimes n})\simeq\mathcal L_{C/S}^{\otimes n}
\]

On a une section \unsure{évidente} de $\underline L$, savoir la section constante $1$ de $\underline O_C$,
\[
\sigma:\underline O_C\longrightarrow\underline L\qquad \sigma\in\Gamma(C,\underline L)
\]
qui définit un homomorphisme (\unsure{d'équations linéaires})
\diagram{\begin{tikzcd}[column sep=small]
\cdots\arrow[r]&\underline L^{\otimes n}\arrow[r,dash]&\underline L^{\otimes(n+1)}\arrow[r]&\cdots\arrow[r]&\underline L^{\otimes(n+\unsure{h})}\arrow[r]&\cdots
\end{tikzcd}}
\margin{$\subset$ faisc. des fonctions rationnelles \unsure{relatives} sur $C/S$}
\unsure{d'où} \deleted{\unsure{un}} des homomorphismes
\diagram{\begin{tikzcd}[column sep=large]
\cdots\arrow[r]&f_*(\underline L^{\otimes n})\arrow[r]&f_*(\underline L^{\otimes(n+1)})\arrow[r]&\cdots
\end{tikzcd}}
i.e.
\diagram{\begin{tikzcd}
\cdots\arrow[r]&\mathcal E_n\arrow[r,dash]&\mathcal E_{n+1}\arrow[r]&\cdots
\end{tikzcd}}
en posant $\mathcal E_n=f_*(\underline L^{\otimes n})$

\mspage{3}
